\documentclass[11pt]{article}
\usepackage[margin=1in]{geometry}
\usepackage{amsmath,amssymb,amsthm}
\usepackage{graphicx}
\usepackage[hidelinks]{hyperref}

\newtheorem{theorem}{Theorem}
\newtheorem{lemma}{Lemma}
\newtheorem{corollary}{Corollary}
\newtheorem{remark}{Remark}

\newcommand{\N}{\mathbb N}
\newcommand{\F}{\mathbb F}
\newcommand{\supp}{\operatorname{supp}}
\newcommand{\spn}{\operatorname{span}}
\newcommand{\Vstar}{V^\ast}
\newcommand{\EVstar}{EV^\ast}
\newcommand{\ellinf}{\ell_\infty}

\title{Free-LCS \(V^\ast\) Variants and the Weak Sequential-Precompact Split}
\author{Run \texttt{fa\_banach\_001}, agent \texttt{agent\_lane\_17}}
\date{June 15, 2026}

\begin{document}
\maketitle

\section*{Status and Source Problem}

This packet gives a consolidated partial answer to Problem 3.18 of Saak
Gabriyelyan's paper \emph{Gelfand--Phillips type properties of locally
convex spaces} \cite{Gab2024GP}. The problem asks:

\begin{center}
\includegraphics[width=\textwidth]{figures/open_problem_crop.png}
\end{center}

The first result below is a full characterization for the properties
\[
  \Vstar_{(p,q)},\qquad \EVstar_{(p,q)},\qquad
  s\Vstar_{(p,q)},\qquad s\EVstar_{(p,q)}
\]
of the free locally convex space \(L(X)\). The second result gives a
positive criterion for the weak sequential-precompact variants
\(ws\Vstar_{(p,q)}\) and \(ws\EVstar_{(p,q)}\), and shows by the example
\(X=[0,1]\) that these weak variants cannot have the same
finite-functionally-bounded characterization as the stronger variants.

\section*{Definitions and Standard Facts Used}

For a Tychonoff space \(X\), write elements of the free locally convex
space \(L(X)\) as finite sums
\[
  \chi=a_1x_1+\cdots+a_nx_n,
\]
where \(x_1,\ldots,x_n\in X\) are distinct and \(a_i\neq 0\). Put
\[
  \|\chi\|_1=|a_1|+\cdots+|a_n|,
  \qquad
  \supp(\chi)=\{x_1,\ldots,x_n\}.
\]
For \(A\subseteq L(X)\), let
\[
  \supp(A)=\bigcup_{\chi\in A}\supp(\chi),\qquad
  C_A=\sup\{\|\chi\|_1:\chi\in A\}.
\]
The dual of \(L(X)\) is \(C(X)\), with pairing
\[
  \langle f,\chi\rangle=a_1f(x_1)+\cdots+a_nf(x_n).
\]

We use the bounded-set criterion for free locally convex spaces recalled
in Theorem 3.17 of \cite{Gab2024GP} from Proposition 2.7 of
\cite{Gab2019LCS}: a subset \(B\subseteq L(X)\) is bounded if and only if
\(\supp(B)\) is functionally bounded in \(X\) and \(C_B<\infty\).

A nonempty set \(A\subseteq E\) is a \((p,q)\)-\(\Vstar\) set if, for
every weakly \(p\)-summable sequence \((\chi_n)\) in the strong dual
\(E'_\beta\), the scalar sequence
\[
  \left(\sup_{a\in A} |\langle \chi_n,a\rangle|\right)_n
\]
belongs to \(\ell_q\), if \(q<\infty\), and to \(c_0\), if
\(q=\infty\). The \((p,q)\)-\(\EVstar\) sets are defined similarly, but
only equicontinuous weakly \(p\)-summable sequences in \(E'_\beta\) are
tested. We use Lemma 7.2 of \cite{Gab2024Pel}: every \((p,q)\)-\(\Vstar\)
set is a \((p,q)\)-\(\EVstar\) set, and every \((p,q)\)-\(\EVstar\) set
is bounded.

\section*{Proof Intuition}

The sufficient direction is finite-dimensional: if all functionally
bounded subsets of \(X\) are finite, then every bounded subset of \(L(X)\)
has finite support and bounded coefficient norm. Since every
\(\EVstar\)-set is bounded, every \(\Vstar\)- or \(\EVstar\)-set is
relatively compact in a finite-dimensional coordinate span.

For the converse, an infinite functionally bounded set \(D=\{x_i\}\)
supports a sequence of finite measures
\[
  a_m=\sum_{i=0}^{m}2^{-i-1}x_i\in L(X).
\]
These partial sums behave like a norm-compact set after restriction to
\(D\), so they are \((p,q)\)-\(\Vstar\) sets for every \(p\le q\). But in
the weak topology of \(L(X)\) they converge pointwise on \(C(X)\) to the
infinite measure \(\sum_i2^{-i-1}\delta_{x_i}\), which is continuous on
the strong dual but is not an element of \(L(X)\). Thus the sequence is a
\(\Vstar\) witness that compactness and sequential compactness fail.

The weak sequential-precompact variants require a different idea. If a
bounded sequence in \(L(X)\) has countable support \(S\), then it is a
uniformly bounded sequence of functionals on the Banach space obtained by
closing the restrictions \(\{f|_S:f\in C(X)\}\) in \(\ell_\infty(S)\).
When that restriction space is separable, its dual ball is weak-star
compact metrizable, and hence the original sequence has a weakly Cauchy
subsequence.

\begin{lemma}[Weighted partial sums are universal \(\Vstar\) sets]
\label{lem:weighted-vstar}
Let \(D=\{x_i:i\in\N\}\) be an infinite functionally bounded subset of a
Tychonoff space \(X\), and put \(w_i=2^{-i-1}\). For
\[
  a_m=\sum_{i=0}^{m} w_i x_i\in L(X),\qquad m\in\N,
\]
the set \(A=\{a_m:m\in\N\}\) is a \((p,q)\)-\(\Vstar\) subset of \(L(X)\)
for every \(1\le p\le q\le\infty\).
\end{lemma}

\begin{proof}
Since \(D\) is functionally bounded, the set \(\{x_i:i\in\N\}\) is a
bounded subset of \(L(X)\). Hence the restriction map
\[
  R:L(X)'_\beta=C(X)_\beta\longrightarrow \ell_\infty,\qquad
  R(f)=(f(x_i))_{i\in\N},
\]
is continuous: its norm is the strong-dual seminorm of uniform convergence
on the bounded set \(\{x_i:i\in\N\}\).

Define a bounded operator \(H:\ell_\infty\to\ell_\infty\) by
\[
  H(y)(m)=\sum_{i=0}^{m}w_i y(i),\qquad y=(y(i))_{i\in\N}.
\]
We first record that \(H\) sends weakly \(p\)-summable sequences in
\(\ell_\infty\) to norm \(\ell_p\)-summable sequences when \(p<\infty\),
and sends weakly null sequences to norm-null sequences when \(p=\infty\).

Let \(1\le p<\infty\), and let \((y_k)\) be weakly \(p\)-summable in
\(\ell_\infty\). The map
\[
  S:(\ell_\infty)'\to \ell_p,\qquad
  S(\varphi)=(\varphi(y_k))_{k\in\N},
\]
is well-defined. Its graph is closed, so by the closed graph theorem
\(S\) is bounded. If \(e_i'\) denotes the \(i\)-th coordinate functional
on \(\ell_\infty\), then
\[
  \big\|(y_k(i))_{k\in\N}\big\|_{\ell_p}
  =\|S(e_i')\|_{\ell_p}\le \|S\|
  \quad\text{for every }i.
\]
Minkowski's inequality gives
\[
\begin{aligned}
  \left(\sum_k \|H(y_k)\|_\infty^p\right)^{1/p}
  &\le
  \left(\sum_k \left(\sum_i w_i |y_k(i)|\right)^p\right)^{1/p}  \\
  &\le
  \sum_i w_i\left(\sum_k |y_k(i)|^p\right)^{1/p}
  \le \|S\|\sum_i w_i <\infty .
\end{aligned}
\]
Thus \((\|H(y_k)\|_\infty)_k\in\ell_p\), hence it belongs to
\(\ell_q\) for every \(q\ge p\).

For \(p=\infty\), the finite-rank truncations
\[
  H_N(y)(m)=\sum_{i=0}^{\min\{m,N\}} w_i y(i)
\]
converge to \(H\) in operator norm, since
\(\|H-H_N\|\le \sum_{i>N}w_i\). Therefore \(H\) is compact. Compact
operators send weakly null sequences in a Banach space to norm-null
sequences, so \(\|H(y_k)\|_\infty\to0\).

Now let \((f_k)\) be weakly \(p\)-summable in \(L(X)'_\beta\). By the
continuity of \(R\), the sequence \(y_k=R(f_k)\) is weakly \(p\)-summable
in \(\ell_\infty\). Moreover,
\[
  \sup_{a\in A}|\langle f_k,a\rangle|
  =
  \sup_m\left|\sum_{i=0}^{m}w_i f_k(x_i)\right|
  =\|H(y_k)\|_\infty .
\]
The preceding paragraph proves that these suprema lie in \(\ell_q\), or
in \(c_0\) when \(q=\infty\). Hence \(A\) is a
\((p,q)\)-\(\Vstar\) set.
\end{proof}

\begin{lemma}[The weighted partial sums are not weakly compact]
\label{lem:not-compact}
With \(D\), \(w_i\), \(a_m\), and \(A\) as in Lemma
\ref{lem:weighted-vstar}, the set \(A\) is neither relatively weakly
compact nor relatively weakly sequentially compact in \(L(X)\).
\end{lemma}

\begin{proof}
Define
\[
  \mu(f)=\sum_{i=0}^{\infty}w_i f(x_i),\qquad f\in C(X).
\]
The series is absolutely convergent because \(D\) is functionally bounded.
Also \(\mu\in L(X)''=(L(X)'_\beta)'\): indeed, the seminorm
\(f\mapsto\sup_i|f(x_i)|\) is continuous on \(L(X)'_\beta\), since
\(\{x_i:i\in\N\}\) is bounded in \(L(X)\), and
\[
  |\mu(f)|\le \left(\sum_iw_i\right)\sup_i|f(x_i)|.
\]
For every \(f\in C(X)\), one has
\[
  \langle f,a_m\rangle=\sum_{i=0}^m w_i f(x_i)\longrightarrow \mu(f).
\]

The functional \(\mu\) is not represented by any element of \(L(X)\). If
\(\mu=\sum_{j=1}^r c_j y_j\) for finitely many \(y_j\in X\), choose
\(N\) with \(x_N\notin \{y_1,\ldots,y_r\}\). Since \(X\) is Tychonoff and
the set \(\{y_1,\ldots,y_r,x_0,\ldots,x_{N-1}\}\) is finite and closed,
there is \(f\in C(X,[0,1])\) such that \(f(x_N)=1\) and \(f\) vanishes on
that finite set. Then the finite sum gives \(\sum_j c_j f(y_j)=0\), while
\[
  \mu(f)=\sum_i w_i f(x_i)\ge w_N>0,
\]
a contradiction.

If \(A\) were relatively weakly compact, its weak closure \(K\) in
\(L(X)\) would be compact. The canonical evaluation map
\[
  J:L(X)_w\to \F^{C(X)},\qquad J(\chi)=(\langle f,\chi\rangle)_{f\in C(X)}
\]
is a topological embedding. The compact set \(J(K)\) is closed in the
Hausdorff product \(\F^{C(X)}\), contains \(J(A)\), and therefore contains
the product limit \((\mu(f))_{f\in C(X)}\) of the sequence \(J(a_m)\).
This would represent \(\mu\) by an element of \(L(X)\), contradicting the
previous paragraph. Hence \(A\) is not relatively weakly compact.

The same pointwise-limit argument shows that no subsequence of
\((a_m)\) weakly converges in \(L(X)\): any subsequential weak limit would
again represent \(\mu\). Therefore \(A\) is not relatively weakly
sequentially compact.
\end{proof}

\begin{theorem}[Characterization for \(V^\ast\), \(EV^\ast\), and their sequential variants]
\label{thm:main}
Let \(1\le p\le q\le\infty\), and let \(X\) be a Tychonoff space. The
following assertions are equivalent:
\begin{enumerate}
\item every functionally bounded subset of \(X\) is finite;
\item \(L(X)\) has the property \(\Vstar_{(p,q)}\);
\item \(L(X)\) has the property \(\EVstar_{(p,q)}\);
\item \(L(X)\) has the property \(s\Vstar_{(p,q)}\);
\item \(L(X)\) has the property \(s\EVstar_{(p,q)}\).
\end{enumerate}
Moreover, if these conditions hold, then \(L(X)\) also has
\(ws\Vstar_{(p,q)}\) and \(ws\EVstar_{(p,q)}\).
\end{theorem}

\begin{proof}
Assume first that every functionally bounded subset of \(X\) is finite.
Let \(B\subseteq L(X)\) be a \((p,q)\)-\(\EVstar\) set. By Lemma 7.2 of
\cite{Gab2024Pel}, \(B\) is bounded. The bounded-set criterion for
\(L(X)\) gives \(C_B<\infty\) and says that \(\supp(B)\) is functionally
bounded; by the present hypothesis, \(F=\supp(B)\) is finite. Hence
\[
  B\subseteq
  \left\{\sum_{x\in F}a_xx:\ \sum_{x\in F}|a_x|\le C_B\right\}.
\]
The set on the right is compact in the finite-dimensional space
\(\spn(F)\). Thus every \((p,q)\)-\(\EVstar\) set is relatively compact,
hence relatively weakly compact, relatively weakly sequentially compact,
and weakly sequentially precompact. Since every \((p,q)\)-\(\Vstar\) set
is a \((p,q)\)-\(\EVstar\) set, all the listed \(V^\ast\), \(EV^\ast\),
sequential, and weak sequential-precompact properties follow.

Conversely, suppose that \(X\) contains an infinite functionally bounded
subset. Lemmas \ref{lem:weighted-vstar} and \ref{lem:not-compact} give a
\((p,q)\)-\(\Vstar\) set \(A\subseteq L(X)\) which is neither relatively
weakly compact nor relatively weakly sequentially compact. Since every
\((p,q)\)-\(\Vstar\) set is a \((p,q)\)-\(\EVstar\) set, the same \(A\)
also witnesses the failure of \(\EVstar_{(p,q)}\) and
\(s\EVstar_{(p,q)}\). Thus none of assertions (2)--(5) can hold.
\end{proof}

\section*{A Positive Criterion for the Weak Variants}

The preceding theorem shows that finite functional boundedness is
sufficient for \(ws\Vstar_{(p,q)}\) and \(ws\EVstar_{(p,q)}\). It is not
necessary. We record a broader sufficient condition which covers all
compact metrizable spaces and gives a concrete separation example.

\begin{theorem}[Restriction separability gives weak sequential precompactness]
\label{thm:restriction-separable}
Let \(X\) be a Tychonoff space. Assume that for every countable
functionally bounded set \(S\subseteq X\), the Banach space
\[
  E_S=\overline{\{f|_S:f\in C(X)\}}^{\|\cdot\|_\infty}
  \subseteq \ell_\infty(S)
\]
is separable. Then every bounded subset of \(L(X)\) is weakly
sequentially precompact.
\end{theorem}

\begin{proof}
Let \(B\subseteq L(X)\) be bounded, and let \((\chi_n)\) be a sequence in
\(B\). By the bounded-set criterion, \(C_B<\infty\) and \(\supp(B)\) is
functionally bounded. Put
\[
  S=\bigcup_{n\in\N}\supp(\chi_n).
\]
Then \(S\) is countable and functionally bounded. By hypothesis,
\[
  E_S=\overline{\{f|_S:f\in C(X)\}}^{\|\cdot\|_\infty}
\]
is a separable Banach space.

Write
\[
  \chi_n=\sum_{x\in F_n}a_{n,x}x,\qquad F_n\subseteq S\ \hbox{finite}.
\]
Each \(\chi_n\) defines a continuous linear functional \(\mu_n\in E_S'\)
by
\[
  \mu_n(g)=\sum_{x\in F_n}a_{n,x}g(x),\qquad g\in E_S.
\]
Moreover
\[
  \|\mu_n\|\le \sum_{x\in F_n}|a_{n,x}|=\|\chi_n\|_1\le C_B.
\]
Thus \((\mu_n)\) is contained in the weak-star compact set
\(C_BB_{E_S'}\). Since \(E_S\) is separable, this compact set is
metrizable for the weak-star topology. Hence \((\mu_n)\) has a weak-star
convergent subsequence, say \((\mu_{n_k})\).

For every \(f\in C(X)\), its restriction \(f|_S\) belongs to \(E_S\), so
the scalar sequence
\[
  \langle f,\chi_{n_k}\rangle
  =
  \sum_{x\in F_{n_k}}a_{n_k,x}f(x)
  =
  \mu_{n_k}(f|_S)
\]
converges. Hence \((\chi_{n_k})\) is weakly Cauchy in \(L(X)\). Since the
original sequence in \(B\) was arbitrary, \(B\) is weakly sequentially
precompact.
\end{proof}

\begin{corollary}[Weak \(V^\ast\) variants under restriction separability]
\label{cor:wsvstar}
Under the hypothesis of Theorem \ref{thm:restriction-separable},
\(L(X)\) has \(ws\Vstar_{(p,q)}\) and \(ws\EVstar_{(p,q)}\) for every
\(1\le p\le q\le\infty\).
\end{corollary}

\begin{proof}
Every \((p,q)\)-\(\Vstar\) set is a \((p,q)\)-\(\EVstar\) set, and every
\((p,q)\)-\(\EVstar\) set is bounded. By Theorem
\ref{thm:restriction-separable}, every bounded subset of \(L(X)\) is
weakly sequentially precompact.
\end{proof}

\begin{corollary}[Compact metrizable spaces and separation]
\label{cor:interval-separation}
If \(K\) is compact metrizable, then \(L(K)\) has
\[
  ws\Vstar_{(p,q)}
  \quad\hbox{and}\quad
  ws\EVstar_{(p,q)}
\]
for every \(1\le p\le q\le\infty\). In particular, \(L([0,1])\) has all
of these weak variants, but fails
\[
  \Vstar_{(p,q)},\quad \EVstar_{(p,q)},\quad
  s\Vstar_{(p,q)},\quad s\EVstar_{(p,q)}
\]
for every \(1\le p\le q\le\infty\).
\end{corollary}

\begin{proof}
Since \(K\) is compact metrizable, \(C(K)\) is separable. For each
countable \(S\subseteq K\), the restriction space
\[
  \overline{\{f|_S:f\in C(K)\}}^{\|\cdot\|_\infty}
\]
is the closure of a continuous image of \(C(K)\), and is therefore
separable. The weak variants follow from Corollary \ref{cor:wsvstar}.

For \(K=[0,1]\), the space \(K\) is infinite and functionally bounded.
Theorem \ref{thm:main} therefore gives the failure of the four stronger
variants listed above.
\end{proof}

\section*{Limitations and Novelty Check}

Theorem \ref{thm:main} characterizes the compact and sequentially compact
\(V^\ast\)-type variants \(V^\ast\), \(EV^\ast\), \(sV^\ast\), and
\(sEV^\ast\). Corollary \ref{cor:interval-separation} shows that the
weaker \(wsV^\ast\) and \(wsEV^\ast\) variants have a different
characterization: finite functional boundedness is sufficient but not
necessary.

A further push on the missing converse for the weak variants did not
produce a justified full characterization. The operator characterization
of \(wsV^\ast\) in \cite{Gab2024Pel} reduces the missing direction to
constructing, from a bad countable functionally bounded support, a
\((q,p)\)-convergent adjoint \(C(X)\to\ell_\infty\) whose associated
sequence in \(L(X)\) has no weakly Cauchy subsequence. The obvious
candidates fail in opposite ways: point masses and block averages over
\(\beta\N\) have no weakly Cauchy subsequences, but coordinate or block
indicator functions are weakly \(1\)-summable and violate the
\(V^\ast\)-test; summable-weight sign families satisfy the \(V^\ast\)
estimates, but diagonal selection and dominated convergence give weakly
Cauchy subsequences.

An earlier packet proved only the sufficient direction for ordinary
\(V^\ast\)-sets, and a later packet recorded only Theorem
\ref{thm:restriction-separable}. The present consolidated packet
supersedes both.

Novelty checks were bounded but targeted. I checked the local run indexes
for arXiv:2406.07178 and the free-lcs \(V^\ast\) keywords, the source
paper \cite{Gab2024GP}, Gabriyelyan's Pe{\l}czy\'nski-set preprint
\cite{Gab2024Pel}, and Gabriyelyan's limited-set preprint
\cite{Gab2024Limited}. I also searched the web for the exact Problem 3.18
wording and close variants such as ``\(V^\ast\) free locally convex
space \(L(X)\)'' and ``Problem 3.18 Gelfand--Phillips type properties''.
No later paper explicitly answering this free-lcs \(V^\ast\) question was
found in that bounded search.

\section*{Human Review Recommendation}

There are two key points to check. First, in Lemma
\ref{lem:weighted-vstar}, verify that the weighted partial-sum operator
\(H:\ell_\infty\to\ell_\infty\) has the asserted convergence property on
weakly \(p\)-summable sequences. Once that is confirmed, the
finite-functional-boundedness characterization follows from standard
free-lcs bounded-set arguments.

Second, in Theorem \ref{thm:restriction-separable}, verify that passing
from a bounded sequence in \(L(X)\) to the weak-star compact metrizable
dual ball of \(E_S\) is legitimate and that weak-star convergence in
\(E_S'\) gives weak Cauchyness in \(L(X)\). This is the entire mechanism
behind the positive weak-variant criterion and the \(X=[0,1]\) separation.

\begin{thebibliography}{9}

\bibitem{Gab2024GP}
S. Gabriyelyan,
\emph{Gelfand--Phillips type properties of locally convex spaces},
arXiv:2406.07178, 2024.

\bibitem{Gab2024Pel}
S. Gabriyelyan,
\emph{Pe{\l}czy{\'n}ski's type sets and Pe{\l}czy{\'n}ski's geometrical
properties of locally convex spaces},
arXiv:2402.08860, 2024.

\bibitem{Gab2024Limited}
S. Gabriyelyan,
\emph{Limited type subsets of locally convex spaces},
arXiv:2403.02016, 2024.

\bibitem{Gab2019LCS}
S. Gabriyelyan,
\emph{Locally convex properties of free locally convex spaces},
J. Math. Anal. Appl. 480 (2019), 123453.

\end{thebibliography}

\end{document}
